\documentclass[12pt]{article}
\usepackage[utf8]{inputenc}
\usepackage[T1]{fontenc}
\usepackage[margin=1in]{geometry}
\usepackage{setspace}
\usepackage{xcolor}
\usepackage[colorlinks=true, urlcolor=blue!60!black, linkcolor=blue!60!black]{hyperref}
\setlength{\parindent}{2em}
\setlength{\parskip}{1.25em}
\renewcommand{\baselinestretch}{1.0}
\begin{document}
\begin{center}
{\LARGE \textit{Griffin v. Breckenridge}, 403 U.S. 88 (1971)\par}
\vspace{0.8em}
{\large Case Brief\par}
\end{center}
\vspace{1em}

\section{Memory Jogger}

Under 42 U.S.C. \S~1985(3), Congress constitutionally reached purely private, racially motivated conspiracies without any requirement of state action, grounded in Section 2 of the Thirteenth Amendment and the federal right of interstate travel, so long as the conspirators acted with racial or class-based invidiously discriminatory animus.

\section{Detailed Case Facts}

On July 2, 1966, petitioners---Eugene Griffin and other Black citizens of the United States residing in Kemper County, Mississippi (which borders Alabama)---were passengers in an automobile owned and driven by R. G. Grady, a resident of Memphis, Tennessee (\href{https://supreme.justia.com/cases/federal/us/403/88/}{403 U.S. at 89--90, 106}; \href{https://www.courtlistener.com/opinion/108362/griffin-v-breckenridge/}{CourtListener}). They were traveling on federal, state, and local highways in and around DeKalb, Mississippi, running errands and visiting friends (\href{https://supreme.justia.com/cases/federal/us/403/88/}{403 U.S. at 90}).

Respondents Lavon Breckenridge and James Calvin Breckenridge, two white adult residents of DeKalb, Mississippi, mistakenly believed that Grady was a civil rights worker for Black Americans (\href{https://supreme.justia.com/cases/federal/us/403/88/}{403 U.S. at 90}; \href{https://www.oyez.org/cases/1970/144}{Oyez}). According to the complaint, the Breckenridges willfully and maliciously conspired to block Grady's car on the public highway, detain its occupants, and assault them with deadly weapons in order to prevent petitioners and other Black Americans from seeking equal protection of the laws and enjoying equal rights, privileges, and immunities---including free speech, assembly, association, movement, personal security, freedom from enslavement, and the right to travel public highways on the same terms as white citizens (\href{https://supreme.justia.com/cases/federal/us/403/88/}{403 U.S. at 90--92}).

Pursuant to this conspiracy, respondents drove their truck into the path of Grady's car to block the road, forced Grady and the passengers out of the vehicle, and prevented their escape by pointing firearms and threatening to kill or injure them (\href{https://supreme.justia.com/cases/federal/us/403/88/}{403 U.S. at 90--91}). While both defendants held the victims at gunpoint, James Calvin Breckenridge clubbed Grady and each of the petitioners on and about the head with a blackjack, pipe, or similar club, inflicting serious physical injuries and severe emotional distress (\href{https://supreme.justia.com/cases/federal/us/403/88/}{403 U.S. at 91}). Neither defendant was a state official or acting under color of state law. Petitioners filed a civil damages action in federal court under 42 U.S.C. \S~1985(3) (originally enacted as \S~2 of the Civil Rights Act of 1871, the Ku Klux Klan Act, 17 Stat. 13) (\href{https://supreme.justia.com/cases/federal/us/403/88/}{403 U.S. at 92, 99}).

\section{Procedural History}

\begin{itemize}
\item \textbf{United States District Court for the Southern District of Mississippi:} The District Court dismissed petitioners' complaint for failure to state a cause of action, relying on \textit{Collins v. Hardyman}, 341 U.S. 651 (1951), which had construed 42 U.S.C. \S~1985(3) to reach only conspiracies under color of state law in order to avoid perceived constitutional difficulties (\href{https://supreme.justia.com/cases/federal/us/403/88/}{403 U.S. at 92}).
\item \textbf{United States Court of Appeals for the Fifth Circuit (1969):} The Fifth Circuit affirmed the dismissal (\textit{Griffin v. Breckenridge}, 410 F.2d 817 (5th Cir. 1969)). Writing for the panel, Judge Goldberg expressed ``serious doubts'' about the ``continued vitality'' of \textit{Collins v. Hardyman} and remarked that it would not surprise the court if \textit{Collins} were disapproved, but concluded that the appellate court was bound to follow existing Supreme Court precedent (\href{https://supreme.justia.com/cases/federal/us/403/88/}{403 U.S. at 92--93}).
\item \textbf{United States Supreme Court (1970--1971):} The Court granted certiorari (397 U.S. 1074), heard oral argument on January 13--14, 1971, and decided the case on June 7, 1971, reversing the Fifth Circuit and remanding to the District Court for further proceedings (\href{https://www.oyez.org/cases/1970/144}{Oyez}; \href{https://supreme.justia.com/cases/federal/us/403/88/}{403 U.S. at 93, 107}).
\end{itemize}

\section{Judicial Votes}

\textbf{Vote:} Unanimous in judgment, 9--0 (8--1 as to Part V-B regarding the right of interstate travel) (\href{https://supreme.justia.com/cases/federal/us/403/88/}{403 U.S. at 89, 107}).

\begin{itemize}
\item \textbf{Majority Opinion:} Justice Potter Stewart, joined in full by Chief Justice Warren E. Burger and Justices Hugo L. Black, William O. Douglas, William J. Brennan, Jr., Byron R. White, Thurgood Marshall, and Harry A. Blackmun, and joined by Justice John M. Harlan II except as to Part V-B.
\item \textbf{Concurrence:} Justice John M. Harlan II filed a brief concurring statement joining the judgment and all of the Court's opinion except Part V-B (\href{https://supreme.justia.com/cases/federal/us/403/88/}{403 U.S. at 107}).
\item \textbf{Dissents:} None.
\end{itemize}

\section{Holding}

\textbf{Judgment:} Reversed and remanded to the United States District Court for the Southern District of Mississippi for further proceedings (\href{https://supreme.justia.com/cases/federal/us/403/88/}{403 U.S. at 107}).

\textbf{Rule Announced:} (1) Statutorily, 42 U.S.C. \S~1985(3) does not require state action and reaches wholly private conspiracies, provided the complaint alleges four elements: defendants (i) conspired or went in disguise on the highway or on another's premises, (ii) for the purpose of depriving any person or class of persons of the equal protection of the laws or equal privileges and immunities under the laws---meaning there must be ``some racial, or perhaps otherwise class-based, invidiously discriminatory animus behind the conspirators' action''---(iii) committed or caused an act in furtherance of the conspiracy, (iv) whereby another was injured in person or property or deprived of having and exercising any right or privilege of a U.S. citizen (\href{https://supreme.justia.com/cases/federal/us/403/88/}{403 U.S. at 102--03}). (2) Constitutionally, Congress possessed authority to reach the racially motivated private highway conspiracy alleged here both under Section 2 of the Thirteenth Amendment (to abolish the ``badges and incidents of slavery'') and under its power to protect the constitutional right of interstate travel against private interference (\href{https://supreme.justia.com/cases/federal/us/403/88/}{403 U.S. at 104--06}).

\textbf{Holding vs. Dicta:} The \textit{holding} is that \S~1985(3) covers purely private conspiracies motivated by racial animus, that petitioners' complaint stated a valid cause of action under \S~1985(3), and that applying the statute to these allegations is a constitutional exercise of Congress's powers under Section 2 of the Thirteenth Amendment and its power to protect interstate travel (\href{https://supreme.justia.com/cases/federal/us/403/88/}{403 U.S. at 101--06}). Two important points are explicitly reserved or constitute \textit{dicta}: first, the Court expressly declined to decide ``whether a conspiracy motivated by invidiously discriminatory intent other than racial bias would be actionable'' under \S~1985(3), making the famous phrase ``or perhaps otherwise class-based'' non-binding \textit{dicta} (\href{https://supreme.justia.com/cases/federal/us/403/88/}{403 U.S. at 102 \& n.9}); second, the Court expressly declined to decide whether Section 5 of the Fourteenth Amendment independently authorizes Congress to reach private conspiracies, leaving that question open (\href{https://supreme.justia.com/cases/federal/us/403/88/}{403 U.S. at 107}; \href{https://constitution.congress.gov/browse/essay/amdt14-S5-2/ALDE_00000851/}{Constitution Annotated}).

\section{Analysis of Opinions}

\subsection{Majority Opinion (Justice Stewart)}

Justice Stewart dismantled the state-action gloss of \textit{Collins v. Hardyman} through four mutually reinforcing statutory arguments. First, the \textit{text} of \S~1985(3) speaks simply of ``two or more persons'' who ``conspire or go in disguise on the highway or on the premises of another''---conduct classically associated with private Ku Klux Klan marauders rather than state officials---and omits any ``under color of state law'' requirement (\href{https://supreme.justia.com/cases/federal/us/403/88/}{403 U.S. at 96--97}). Second, \textit{precedent} construing Reconstruction civil rights statutes accorded them ``a sweep as broad as [their] language,'' and cases interpreting criminal analogues such as 18 U.S.C. \S~241 (\textit{United States v. Williams}, 341 U.S. 70 (1951)) and Rev. Stat. \S~5519 (\textit{United States v. Harris}, 106 U.S. 629 (1883)) recognized their reach over private conduct (\href{https://supreme.justia.com/cases/federal/us/403/88/}{403 U.S. at 97--98}). Third, \textit{structural context} within the Civil Rights Act of 1871 showed Congress explicitly addressed state-action deprivations in \S~1 (now 42 U.S.C. \S~1983), interference with state officials in a separate clause of \S~1985(3), and massive lawlessness overriding state authorities in \S~3; reading a state-action limit into the clause at issue would render it redundant (\href{https://supreme.justia.com/cases/federal/us/403/88/}{403 U.S. at 98--99}). Fourth, \textit{legislative history}---including statements by Representatives Shellabarger, Willard, Shanks, and Coburn and Senator Pool---confirmed Congress intended to reach private conspirators directly (\href{https://supreme.justia.com/cases/federal/us/403/88/}{403 U.S. at 99--101}).

To avoid turning \S~1985(3) into a ``general federal tort law'' and hitting ``constitutional shoals,'' Stewart read the statutory phrase ``for the purpose of depriving ... of the equal protection of the laws, or of equal privileges and immunities under the laws'' as requiring ``some racial, or perhaps otherwise class-based, invidiously discriminatory animus'' (\href{https://supreme.justia.com/cases/federal/us/403/88/}{403 U.S. at 101--02}). Applying the four statutory elements, petitioners' allegations of a racially motivated ambush and beating on a public highway easily stated a claim (\href{https://supreme.justia.com/cases/federal/us/403/88/}{403 U.S. at 102--03}).

Turning to constitutional authority, Stewart first discarded the outdated non-severability rule of \textit{United States v. Harris} and \textit{Baldwin v. Franks}, holding under \textit{United States v. Raines}, 362 U.S. 17 (1960), that \S~1985(3) need only be constitutional as applied to the complaint (\href{https://supreme.justia.com/cases/federal/us/403/88/}{403 U.S. at 104}). He then identified two independent sources of congressional power that require no state action. In Part V-A, relying on the \textit{Civil Rights Cases}, 109 U.S. 3 (1883), and \textit{Jones v. Alfred H. Mayer Co.}, 392 U.S. 409 (1968), he held that Section 2 of the Thirteenth Amendment empowers Congress to rationally determine the ``badges and incidents of slavery'' and create a civil remedy for Black citizens subjected to racially motivated private conspiracies (\href{https://supreme.justia.com/cases/federal/us/403/88/}{403 U.S. at 104--05}; \href{https://law.justia.com/constitution/us/amendment-13/01-abolition-of-slavery.html}{Justia Constitution Annotated}). In Part V-B, relying on \textit{United States v. Guest}, 383 U.S. 745 (1966), and \textit{Shapiro v. Thompson}, 394 U.S. 618 (1969), he held that the constitutional right of interstate travel is assertable against private interference, and petitioners could prove at trial that the Kemper County highway attack targeted interstate travel or association with out-of-state civil rights workers (\href{https://supreme.justia.com/cases/federal/us/403/88/}{403 U.S. at 105--06}). In Part V-C, the Court expressly reserved whether Section 5 of the Fourteenth Amendment also supported the statute (\href{https://supreme.justia.com/cases/federal/us/403/88/}{403 U.S. at 107}).

\subsection{Concurrence (Justice Harlan)}

Justice Harlan joined the judgment and the entire opinion except Part V-B. In a concise two-sentence concurrence, he wrote that he found it ``unnecessary to rely on the `right of interstate travel' as a premise for justifying federal jurisdiction under \S~1985(3)'' (\href{https://supreme.justia.com/cases/federal/us/403/88/}{403 U.S. at 107}). Because Section 2 of the Thirteenth Amendment fully sustained Congress's power to reach the racially motivated private conspiracy alleged in the complaint, Harlan saw no need to invoke the doctrinally contested, unenumerated right of interstate travel as an alternative constitutional basis.

\subsection{Dissents}

No dissenting opinions were filed.

\section{Examples: Future Applications}

\begin{enumerate}
\item \textbf{Same-Side Application (Private Racially Motivated Interference With Interstate Travel and Property):} Members of a private white supremacist group conspire to vandalize the home and slash the tires of a Black family that recently moved across state lines into a predominantly white town, intending to drive them out of the neighborhood because of their race. \textit{Why Same-Side:} Like \textit{Griffin}, this is a purely private conspiracy driven by overt racial animus aiming to deprive Black citizens of basic rights secured to all free persons and to impair interstate movement, squarely within \S~1985(3) and Congress's Thirteenth Amendment and interstate-travel powers (\href{https://supreme.justia.com/cases/federal/us/403/88/}{403 U.S. at 102--06}).

\item \textbf{Same-Side Application (Private Conspiracy Targeting Civil Rights Volunteers for Assisting Black Voters):} Two private citizens conspire to ambush and beat an out-of-state white volunteer who is driving Black residents to a voter-registration rally, intending to intimidate Black residents from participating in civic life. \textit{Why Same-Side:} Even though the direct target is not Black (just as R. G. Grady was a non-party believed to be a civil rights worker), the conspiracy is driven by racial animus aimed at depriving Black citizens of equal rights and deterring association with out-of-state civil rights workers (\href{https://supreme.justia.com/cases/federal/us/403/88/}{403 U.S. at 103, 106}).

\item \textbf{Opposite-Side Application (Private Economic/Labor Conspiracy Against Nonunion Workers):} A group of union supporters conspires to assault nonunion construction workers and destroy equipment at a private job site to force the contractor to hire union labor. \textit{Why Opposite-Side:} In \href{https://caselaw.findlaw.com/court/us-supreme-court/463/825.html}{\textit{United Brotherhood of Carpenters \& Joiners v. Scott}}, 463 U.S. 825 (1983), the Supreme Court held that \S~1985(3) does not reach wholly private conspiracies motivated by economic or commercial animus, and that a private conspiracy to infringe First Amendment associational rights is not actionable under \S~1985(3) without state involvement (\href{https://caselaw.findlaw.com/court/us-supreme-court/463/825.html}{463 U.S. at 833, 837--38}).

\item \textbf{Opposite-Side Application (Private Protesters Blocking Access to Abortion Clinics):} Anti-abortion demonstrators conspire to blockade entrances to women's health clinics to stop abortions from being performed, incidentally impeding patients who traveled from neighboring states. \textit{Why Opposite-Side:} In \href{https://supreme.justia.com/cases/federal/us/506/263/}{\textit{Bray v. Alexandria Women's Health Clinic}}, 506 U.S. 263 (1993), the Court held that opposition to abortion does not constitute ``class-based, invidiously discriminatory animus'' against women as a class under \textit{Griffin}, and that an incidental effect on interstate travelers does not satisfy the requirement that the conspiracy be ``aimed at'' a right protected against private encroachment (\href{https://supreme.justia.com/cases/federal/us/506/263/}{506 U.S. at 269--70, 274--75}).

\item \textbf{Fence-Sitter Application (Private Conspiracy Targeting a Religious Minority on a Public Highway):} Members of a private militia conspire to block cars driven by visibly identifiable members of a religious minority on an interstate highway near a state border, assaulting them to drive their congregation out of the state. \textit{Why Fence-Sitter:} Unlike \textit{Carpenters} or \textit{Bray}, the conspirators act with overt group-based animus and directly target interstate travel (a right protected against private interference under Part V-B of \textit{Griffin}). However, because \textit{Griffin} expressly reserved whether non-racial class-based animus is actionable (\href{https://supreme.justia.com/cases/federal/us/403/88/}{403 U.S. at 102 n.9}) and subsequent courts have split over whether religious groups qualify as a protected class under \S~1985(3) and whether Congress's Thirteenth Amendment power extends beyond race or ancestry, the outcome turns on whether the court treats religious animus as falling within \S~1985(3)'s ``otherwise class-based'' scope (\href{https://constitution.congress.gov/browse/essay/amdt14-S5-2/ALDE_00000851/}{Constitution Annotated}; \href{https://repository.law.umich.edu/cgi/viewcontent.cgi?article=2822&context=mlr}{Michigan Law Review}).
\end{enumerate}

\section{Critique}

\textbf{Originalist and Textualist Critiques:} Originalist and textualist commentators raise two principal objections to \textit{Griffin}. First, some critics argue that the Court's reliance on \textit{Jones v. Alfred H. Mayer Co.} to uphold \S~1985(3) under Section 2 of the Thirteenth Amendment stretches the ``badges and incidents of slavery'' far beyond the Amendment's original public meaning, which they contend was limited to abolishing chattel slavery, involuntary servitude, and legal disabilities approximating enslavement rather than granting Congress broad police power over private torts (\href{https://scholars.law.unlv.edu/cgi/viewcontent.cgi?article=1013&context=amend13_bi}{UNLV Scholarly Commons}; \href{https://scholarship.shu.edu/cgi/viewcontent.cgi?article=1935&context=shlr}{Seton Hall Law Review}). Second, textualists note an internal tension in Stewart's statutory construction: after insisting that \S~1985(3)'s text must be given its literal sweep to cover private conspiracies, the Court simultaneously divorced the phrase ``equal protection of the laws'' from its Fourteenth Amendment state-action framework and grafted onto it an atextual ``class-based, invidiously discriminatory animus'' requirement to prevent the statute from becoming a ``general federal tort law'' (\href{https://scholarship.shu.edu/cgi/viewcontent.cgi?article=1935&context=shlr}{Seton Hall Law Review}; \href{https://repository.law.umich.edu/cgi/viewcontent.cgi?article=2822&context=mlr}{Michigan Law Review}).

\textbf{Progressive and Doctrinal Perspectives:} Progressive and civil rights scholars applaud \textit{Griffin} for resurrecting the Ku Klux Klan Act from a century of judicial narrowing and recognizing that the Thirteenth Amendment provides a robust, state-action-free foundation for combating private racial terror (\href{https://lawreview.law.ucdavis.edu/sites/g/files/dgvnsk15026/files/media/documents/DavisVol40No4_Carter.pdf}{UC Davis Law Review}; \href{https://scholarship.law.bu.edu/cgi/viewcontent.cgi?article=3274&context=faculty_scholarship}{Boston University Law}). Yet progressive critics lament that by avoiding Section 5 of the Fourteenth Amendment and inventing the ambiguous ``racial, or perhaps otherwise class-based, invidiously discriminatory animus'' gatekeeper, \textit{Griffin} planted the seeds for the Supreme Court's later retrenchment in \href{https://caselaw.findlaw.com/court/us-supreme-court/463/825.html}{\textit{Carpenters v. Scott}} (1983) and \href{https://supreme.justia.com/cases/federal/us/506/263/}{\textit{Bray v. Alexandria Women's Health Clinic}} (1993), which severely curtailed \S~1985(3)'s utility for non-racial civil rights plaintiffs and left lower courts mired in doctrinal confusion over which classes and underlying rights qualify (\href{https://constitution.congress.gov/browse/essay/amdt14-S5-2/ALDE_00000851/}{Constitution Annotated}; \href{https://repository.law.umich.edu/cgi/viewcontent.cgi?article=2822&context=mlr}{Michigan Law Review}).

\textbf{Underlying Competing Values:} \textit{Griffin} exposes the enduring clash between \textit{federalism and enumerated powers} (preserving state primacy over ordinary private tort and criminal law and honoring the Fourteenth Amendment's state-action boundary) and \textit{national civil rights enforcement} (empowering Congress to eradicate private racial violence and intimidation when state institutions fail or decline to protect vulnerable citizens).

\section{Key Quotations}

\begin{quote}
``It is thus evident that all indicators -- text, companion provisions, and legislative history -- point unwaveringly to \S~1985(3)'s coverage of private conspiracies.''
\end{quote}
\noindent --- Justice Potter Stewart, writing for the Court, \href{https://supreme.justia.com/cases/federal/us/403/88/}{403 U.S. at 101} (\href{https://www.courtlistener.com/opinion/108362/griffin-v-breckenridge/}{CourtListener}).

\begin{quote}
``The language requiring intent to deprive of equal protection, or equal privileges and immunities, means that there must be some racial, or perhaps otherwise class-based, invidiously discriminatory animus behind the conspirators' action. The conspiracy, in other words, must aim at a deprivation of the equal enjoyment of rights secured by the law to all.''
\end{quote}
\noindent --- Justice Potter Stewart, \href{https://supreme.justia.com/cases/federal/us/403/88/}{403 U.S. at 102} (\href{https://www.courtlistener.com/opinion/108362/griffin-v-breckenridge/}{CourtListener}).

\begin{quote}
``We can only conclude that Congress was wholly within its powers under \S~2 of the Thirteenth Amendment in creating a statutory cause of action for Negro citizens who have been the victims of conspiratorial, racially discriminatory private action aimed at depriving them of the basic rights that the law secures to all free men.''
\end{quote}
\noindent --- Justice Potter Stewart, \href{https://supreme.justia.com/cases/federal/us/403/88/}{403 U.S. at 105} (\href{https://www.courtlistener.com/opinion/108362/griffin-v-breckenridge/}{CourtListener}).

\begin{quote}
``Our cases have firmly established that the right of interstate travel is constitutionally protected, does not necessarily rest on the Fourteenth Amendment, and is assertable against private, as well as governmental, interference.''
\end{quote}
\noindent --- Justice Potter Stewart (Part V-B, joined by 8 Justices), \href{https://supreme.justia.com/cases/federal/us/403/88/}{403 U.S. at 105--06} (\href{https://www.courtlistener.com/opinion/108362/griffin-v-breckenridge/}{CourtListener}).

\begin{quote}
``I agree with the Court's opinion, except that I find it unnecessary to rely on the `right of interstate travel' as a premise for justifying federal jurisdiction under \S~1985(3). With that reservation, I join the opinion and judgment of the Court.''
\end{quote}
\noindent --- Justice John M. Harlan II, concurring, \href{https://supreme.justia.com/cases/federal/us/403/88/}{403 U.S. at 107} (\href{https://www.courtlistener.com/opinion/108362/griffin-v-breckenridge/}{CourtListener}).

\end{document}
